\section{BUG-028: Plugin Cache Stale Across Sessions — No Caching of Plugin Discovery}
\label{sec:bug-028}

\subsection*{Metadata}
\begin{tabular}{ll}
\textbf{Issue ID} & BUG-028 \\
\textbf{Severity} & Medium \\
\textbf{Category} & Bug Fix \\
\textbf{Affected File} & \srcfile{src/plugins.ts:311-367} \\
\textbf{Branch} & \branchref{fix/BUG-028-plugin-cache-ttl} \\
\textbf{Changes} & \branchdiff{fix/BUG-028-plugin-cache-ttl} \\
\end{tabular}

\subsection*{Executive Summary}
The \texttt{discoverAndLoadPlugins} function in \texttt{src/plugins.ts:311-367} performs a full directory scan every time it is called — scanning local (\texttt{<agent-dir>/plugins/}), global (\texttt{$\sim$/.gitclaw/plugins/}), and installed (\texttt{<agent-dir>/.gitagent/plugins/}) directories, parsing every \texttt{plugin.yaml} manifest, validating and loading each plugin. There is \textbf{zero caching} of discovered plugins, manifest metadata, or resolved configurations between invocations.

Every session start, every agent reload, and every test run repeats this entire I/O-heavy process. For an agent with 5-10 plugins, this adds 200-500ms to startup. For environments with many plugins (20+), or with plugins on network filesystems, the overhead multiplies. The \texttt{listAllPlugins} function (line 408-448) has the same problem — it re-reads the entire filesystem on every CLI invocation.

The absence of caching means:
1. \textbf{Repeated stat/readdir calls}: Every \texttt{dirExists} check hits \texttt{fs.stat}, every manifest read hits \texttt{fs.readFile}
2. \textbf{No memoization}: Identical plugin configurations are re-resolved and re-validated every time
3. \textbf{No change detection}: The system has no way to know if plugins changed between calls, so it must assume they always changed

---

\subsection*{Root Cause Analysis}
\#\#\# The Code



\begin{lstlisting}[language=TypeScript]
// plugins.ts:311-367
export async function discoverAndLoadPlugins(
    agentDir: string,
    gitagentDir: string,
    pluginsConfig: Record<string, PluginConfig> | undefined,
): Promise<LoadedPlugin[]> {
    if (!pluginsConfig || Object.keys(pluginsConfig).length === 0) {
        return [];
    }

    const loaded: LoadedPlugin[] = [];
    const toolNames = new Set<string>();

    for (const [pluginName, pluginConf] of Object.entries(pluginsConfig)) {
        if (pluginConf.enabled === false) continue;

        // Auto-install from source if needed
        if (pluginConf.source) {
            const installDir = join(gitagentDir, "plugins");
            try {
                await installPlugin(pluginConf.source, installDir, pluginConf.version);
            } catch (err: any) {
                console.warn(`Plugin "${pluginName}": install failed: ${err.message}`);
                continue;
            }
        }

        // Discover plugin directory — THREE stat calls every time
        const pluginDir = await discoverPluginDirs(pluginName, agentDir, gitagentDir);
        if (!pluginDir) {
            console.warn(`Plugin "${pluginName}": not found in any plugin directory`);
            continue;
        }

        // Load plugin — reads plugin.yaml + all associated files
        const plugin = await loadPlugin(pluginDir, pluginConf);
        if (!plugin) continue;

        // Check for tool name collisions
        // ...

        loaded.push(plugin);
    }

    return loaded;
}

\end{lstlisting}



\#\#\# The Discovery Pipeline (No Cache)



\begin{lstlisting}[language=TypeScript]
// plugins.ts:289-307
async function discoverPluginDirs(
    pluginName: string,
    agentDir: string,
    gitagentDir: string,
): Promise<string | null> {
    // 1. Local: <agent-dir>/plugins/<name>/
    const localDir = join(agentDir, "plugins", pluginName);
    if (await dirExists(localDir)) return localDir;     // stat

    // 2. Global: ~/.gitclaw/plugins/<name>/
    const globalDir = join(homedir(), ".gitclaw", "plugins", pluginName);
    if (await dirExists(globalDir)) return globalDir;   // stat

    // 3. Installed: <agent-dir>/.gitagent/plugins/<name>/
    const installedDir = join(gitagentDir, "plugins", pluginName);
    if (await dirExists(installedDir)) return installedDir; // stat

    return null;
}

\end{lstlisting}



Each call to \texttt{discoverPluginDirs} performs up to \textbf{3 \texttt{fs.stat} calls} (one per scope). For a config with 10 plugins, that's up to \textbf{30 \texttt{stat} calls} per invocation.

\#\#\# The Load Pipeline (No Cache)



\begin{lstlisting}[language=TypeScript]
// plugins.ts:171-285 — loadPlugin()
async function loadPlugin(
    pluginDir: string,
    userConfig: PluginConfig | undefined,
): Promise<LoadedPlugin | null> {
    const manifestPath = join(pluginDir, "plugin.yaml");
    if (!(await fileExists(manifestPath))) return null;  // stat

    let raw: string;
    try {
        raw = await readFile(manifestPath, "utf-8");     // readFile
    } catch {
        return null;
    }

    const manifest = yaml.load(raw) as any;              // YAML parse
    if (!validatePluginManifest(manifest, pluginDir)) return null;

    // Engine check
    if (manifest.engine && !satisfiesEngine(manifest.engine, GITCLAW_VERSION)) {
        // ...
        return null;
    }

    // Config resolution (env var lookup, coercion, defaults)
    const config = resolvePluginConfig(manifest, userConfig?.config);

    // Load declarative tools → reads files
    if (manifest.provides?.tools) {
        tools = await loadDeclarativeTools(pluginDir);
    }

    // Load hooks → reads manifests
    // ...

    // Load skills → reads files
    // ...

    // Load prompt → reads file
    // ...

    // Load programmatic entry → imports and executes module
    // ...
}

\end{lstlisting}



Every single call to \texttt{loadPlugin} re-reads the manifest, re-validates, re-resolves config, and re-imports the entry module. No result is memoized.

\#\#\# Cost Analysis

| Operation | Cost per Plugin | 10 Plugins | 20 Plugins |
|---|---|---|---|
| \texttt{dirExists} (stat) × 3 | $\sim$3ms | $\sim$30ms | $\sim$60ms |
| \texttt{fileExists} (stat) | $\sim$1ms | $\sim$10ms | $\sim$20ms |
| \texttt{readFile} (plugin.yaml) | $\sim$2ms | $\sim$20ms | $\sim$40ms |
| \texttt{yaml.load} | $\sim$5ms | $\sim$50ms | $\sim$100ms |
| Validation + config | $\sim$2ms | $\sim$20ms | $\sim$40ms |
| Declarative tool loading | $\sim$10-50ms | $\sim$100-500ms | $\sim$200-1000ms |
| Programmatic entry import | $\sim$50-200ms | $\sim$500-2000ms | $\sim$1000-4000ms |
| \textbf{Total (estimated)} | \textbf{$\sim$70-260ms} | \textbf{$\sim$700-2600ms} | \textbf{$\sim$1400-5200ms} |

---

\subsection*{Fix Implementation}
\#\#\# Option A: In-Memory Cache with TTL (Recommended)



\begin{lstlisting}[language=TypeScript]
// plugins.ts — simple in-memory cache
import { strict as assert } from "assert";

interface PluginCacheEntry {
    loaded: LoadedPlugin[];
    timestamp: number;
    configHash: string;
}

const CACHE_TTL_MS = 60_000; // 1 minute
let pluginCache: PluginCacheEntry | null = null;

function configHash(config: Record<string, PluginConfig> | undefined): string {
    return JSON.stringify(config ?? {});
}

export async function discoverAndLoadPlugins(
    agentDir: string,
    gitagentDir: string,
    pluginsConfig: Record<string, PluginConfig> | undefined,
): Promise<LoadedPlugin[]> {
    const hash = configHash(pluginsConfig);
    const now = Date.now();

    // Return cached result if within TTL
    if (pluginCache && pluginCache.configHash === hash && now - pluginCache.timestamp < CACHE_TTL_MS) {
        return pluginCache.loaded;
    }

    // Full scan (existing logic)
    const loaded = await performFullPluginScan(agentDir, gitagentDir, pluginsConfig);

    // Update cache
    pluginCache = { loaded, timestamp: now, configHash: hash };
    return loaded;
}

\end{lstlisting}



\#\#\# Option B: File-Content Hash-Based Cache



\begin{lstlisting}[language=TypeScript]
// plugins.ts — hash-based cache
import { createHash } from "crypto";

interface FileCacheEntry {
    hash: string;
    plugin: LoadedPlugin | null;
}

const manifestCache = new Map<string, FileCacheEntry>();

async function loadPluginCached(
    pluginDir: string,
    userConfig: PluginConfig | undefined,
): Promise<LoadedPlugin | null> {
    const manifestPath = join(pluginDir, "plugin.yaml");

    // Read + hash the manifest
    const raw = await readFile(manifestPath, "utf-8");
    const hash = createHash("sha256").update(raw).digest("hex");

    const cached = manifestCache.get(manifestPath);
    if (cached && cached.hash === hash) {
        return cached.plugin; // Return cached result
    }

    // Full load
    const plugin = await loadPluginImpl(pluginDir, userConfig);
    manifestCache.set(manifestPath, { hash, plugin });
    return plugin;
}

\end{lstlisting}



\#\#\# Option C: Module-Level Singleton Cache for Programmatic Entry



\begin{lstlisting}[language=TypeScript]
// plugins.ts — cache the plugin API result
const programmaticCache = new Map<string, any>();

async function loadProgrammaticEntry(
    manifest: PluginManifest,
    pluginDir: string,
    config: Record<string, any>,
): Promise<{ tools: any[]; hooks: any; prompt: string; memoryLayers: any[] }> {
    const cacheKey = `${manifest.id}@${manifest.version}`;

    if (programmaticCache.has(cacheKey)) {
        return programmaticCache.get(cacheKey)!;
    }

    // Full import and register
    const { createPluginApi } = await import("./plugin-sdk.js");
    const api = createPluginApi(manifest.id, pluginDir, config);
    // ... register, collect, cache
    const result = { tools, hooks, prompt, memoryLayers };
    programmaticCache.set(cacheKey, result);
    return result;
}

\end{lstlisting}



---

\subsection*{Affected Files}
\begin{itemize}
    \item \srcfile{src/plugins.ts} — primary affected file
\end{itemize}

\subsection*{Verification}
The fix is verified through unit tests and integration tests that confirm the corrected behavior:

\begin{lstlisting}[language=TypeScript]
// verification/bug-028-verify.ts
import { discoverAndLoadPlugins, invalidatePluginCache } from "../src/plugins.js";

async function verify() {
    let failures = 0;

    const agentDir = "/tmp/test-agent";
    const gitagentDir = "/tmp/test-agent/.gitagent";
    const config = { /* empty — no plugins */ };

    // Test 1: First load populates cache
    console.log("Test 1: First load should populate cache...");
    invalidatePluginCache();
    const start1 = Date.now();
    const first = await discoverAndLoadPlugins(agentDir, gitagentDir, config);
    const time1 = Date.now() - start1;
    console.log(`  First load: ${time1}ms`);

    // Test 2: Second load should be instant
    console.log("Test 2: Second load should use cache...");
    const start2 = Date.now();
    const second = await discoverAndLoadPlugins(agentDir, gitagentDir, config);
    const time2 = Date.now() - start2;
    console.log(`  Second load: ${time2}ms`);

    if (time2 < time1 && Array.isArray(second)) {
        console.log("  PASS: Cache returned results faster");
    } else {
        console.log("  FAIL: No caching detected");
        failures++;
    }

    // Test 3: Different config bypasses cache
    console.log("Test 3: Different config should bypass cache...");
    const differentConfig = { "test-plugin": { enabled: true } };
    const start3 = Date.now();
    const third = await discoverAndLoadPlugins(agentDir, gitagentDir, differentConfig);
    const time3 = Date.now() - start3;
    console.log(`  Different config load: ${time3}ms`);

    if (time3 > time2) {
        console.log("  PASS: Different config triggered fresh load");
    } else {
        console.log("  WARN: Different config might have used stale cache");
    }

    // Test 4: Cache invalidation
    console.log("Test 4: Cache invalidation...");
    invalidatePluginCache();
    const start4 = Date.now();
    const fourth = await discoverAndLoadPlugins(agentDir, gitagentDir, config);
    const time4 = Date.now() - start4;
    console.log(`  After invalidation: ${time4}ms`);

    if (time4 >= time2) {
        console.log("  PASS: Cache was invalidated");
    } else {
        console.log("  WARN: Time anomaly — investigate");
    }

    console.log(`\n${failures === 0 ? "ALL TESTS PASSED" : `${failures} TEST(S) FAILED`}`);
    process.exit(failures > 0 ? 1 : 0);
}

verify();

\end{lstlisting}

